% $Id: INF_Diplomarbeit_example.tex 2733 2010-12-19 07:15:44Z tkren $
%
% TU Wien - Faculty of Informatics
% thesis template
%
% This template is using the memoir document class, see
% <http://www.ctan.org/tex-archive/macros/latex/contrib/memoir/memman.pdf>
% <http://www.ctan.org/info/latex-samples/MemoirChapStyles/MemoirChapStyles.pdf>
%
% For questions and comments send an email to
% Thomas Krennwallner <tkren@kr.tuwien.ac.at>
%

								
\documentclass[a4paper,11pt]{memoir}
\chapterstyle{veelo}

\setcounter{secnumdepth}{4}
\setcounter{tocdepth}{4}
\setsecnumdepth{subsubsection}% alternative B

\usepackage{standalone}
\usepackage{HTBLA}
\usepackage{multirow}
\usepackage{array}
\usepackage{color}
\usepackage{colortbl}
\usepackage{mathtools}
\usepackage{subcaption} % for subfigures
\usepackage{caption}
\usepackage{xargs}
\usepackage[ngerman,colorinlistoftodos]{todonotes}
\usepackage{lipsum}  

\newcommandx{\unsure}[2][1=]{\todo[linecolor=red,backgroundcolor=red!25,bordercolor=red,#1]{#2}}
\newcommandx{\change}[2][1=]{\todo[linecolor=blue,backgroundcolor=blue!25,bordercolor=blue,#1]{#2}}
\newcommandx{\info}[2][1=]{\todo[linecolor=OliveGreen,backgroundcolor=OliveGreen!25,bordercolor=OliveGreen,#1]{#2}}
\newcommandx{\improvement}[2][1=]{\todo[linecolor=Plum,backgroundcolor=Plum!25,bordercolor=Plum,#1]{#2}}
\newcommandx{\thiswillnotshow}[2][1=]{\todo[disable,#1]{#2}}


\usepackage{longtable}
%\usepackage{minted}
\usepackage[boxed,german]{algorithm2e}
\usepackage{footnote}
\usepackage{listings}

\makesavenoteenv{tabular}
\makesavenoteenv{table}
\usepackage{lscape}

% nice tables
\usepackage{booktabs}

% colored tables
\usepackage{colortbl}

% TikZ ist kein Zeichenprogramm
\usepackage{tikz}
\usepackage{tikz-timing}
\usetikzlibrary{shapes.geometric,arrows,shapes,snakes,automata,backgrounds,petri,circuits.logic.IEC,circuits.logic.US,circuits.ee.IEC}
\usepackage[europeanresistors,americaninductors]{circuitikz}

\usepackage{lipsum}  



%%%%.........


\thesistitle{Titel der Arbeit}
\selectlanguage{ngerman}
\thesisdate{\today{}}
\thesislocation{Graz}

%\thesisdegree{Ingenieur}
\thesiscurriculum{Elektronik und technische Informatik}
\thesisverfassung{Verfasser}

\thesisauthoreins{Author Eins}
\thesisauthorzwei{Author Zwei}

\thesisauthoraddresseins{Adresse 1, AT-8051 Graz}
\thesisauthoraddresszwei{Adresse 2, AT-8051 Graz}

\thesismatrikelnoeins{0123456/789}
\thesismatrikelnozwei{0123456/789}

\thesisbetreuung{Betreuer}
\thesisbetreins{Dipl.-Ing. Christian Trenkwalder, BEd}
%\thesisbetrzwei{Univ.-ass. Dipl.-Ing. Robert Navjvirt}

\usepackage{prettyref}

\bibliographystyle{alpha}

% run makeglossaries after the compilation
\usepackage[toc]{glossaries} % nomain, if you define glossaries in a file, and you use \include{INP-00-glossary}
\makeglossaries
%\usepackage[xindy]{imakeidx}

% run makeindex after the compilation
\usepackage{makeidx}
%\makeindex

\makeindex

%\newglossaryentry{citation}{name={Citation},%
%  description={a reference to a publication. (See also \gls{reference})}%
%}
%
%\newglossaryentry{reference}{name={Reference},%
%  description={citing a publication. (See also \gls{citation})}%
%}
%

%\newglossaryentry{api}
%{
%    name={API},
%    description={An Application Programming Interface (API) is a particular set
%    of rules and specifications that a software program can follow to access and make use of the services and resources provided by another particular software program that implements that API},
%    first={Application Programming Interface (API)},
%    long={Application Programming Interface}
%}


% should be the last package
\usepackage{hyperref}
\usepackage[figure]{hypcap}
\usepackage{graphicx}

\hypersetup
{%
pdftitle = {Titel},
pdfsubject = {Diplomarbeit},
pdfauthor = {Author1,Author2},
colorlinks = {true}
} 



\begin{document}
%\renewcommand\listingscaption{Auflistung}

\nouppercaseheads

% the front matter starts
\frontmatter

%
% titlepages
%
% $Id: titlepage.tex 2165 2010-08-07 04:42:52Z tkren $
%
% TU Wien - Faculty of Informatics
% thesis titlepage
%
% This titlepage is using the geometry package, see
% <http://www.ctan.org/macros/latex/contrib/geometry/geometry.pdf>
%
% For questions and comments send an email to
% Thomas Krennwallner <tkren@kr.tuwien.ac.at>
%

% setup page dimensions for titlepage
\newgeometry{left=2.4cm,right=2.4cm,bottom=2.5cm,top=2cm}

% force baselineskip and parindent
\newlength{\tmpbaselineskip}
\setlength{\tmpbaselineskip}{\baselineskip}
\setlength{\baselineskip}{13.6pt}
\newlength{\tmpparindent}
\setlength{\tmpparindent}{\parindent}
\setlength{\parindent}{17pt}

% first titlepage
\thispagestyle{tuinftitlepage}

%
% Kludge: for each titlepage set \pagenumbering to a different
% style. This is used to fix a problem with hyperref, because there
% are multiple "page 1" and hyperref hates that
%
\pagenumbering{Alph}

\begin{center}
%
\begin{minipage}[t][5.2cm][b]{\linewidth}%
    \centering%
    %\renewcommand{\baselinestretch}{0.9} % very long titles need to tweak baselinestretch...
    \thesistitlefontLARGE\sffamily\bfseries\tuinfthesistitle%
\end{minipage}


\vspace{1.3cm}

{\thesistitlefontLARGE\sffamily \tuinfthesistype}

\vspace{6mm}

%{\thesistitlefontlarge\sffamily zur Erlangung des akademischen Grades}

%\vspace{6mm}

%{\thesistitlefontLARGE\sffamily\bfseries \tuinfthesisdegree}

\vspace{6mm}

{\thesistitlefontlarge\sffamily im Rahmen der Fachrichtung}

\vspace{6mm}

{\thesistitlefontLarge\sffamily\bfseries \tuinfthesiscurriculum}

\vspace{6.5mm}

{\thesistitlefontlarge\sffamily eingereicht von}

\vspace{6mm}

{\thesistitlefontLarge\sffamily\bfseries \tuinfthesisauthoreins}

\vspace{1.5mm}

{\thesistitlefontlarge\sffamily Schülerkennzahl \tuinfthesismatrikelnoeins} 

\vspace{6mm}

{\thesistitlefontLarge\sffamily\bfseries \tuinfthesisauthorzwei}

\vspace{1.5mm}

{\thesistitlefontlarge\sffamily Schülerkennzahl \tuinfthesismatrikelnozwei} 

\vspace{1.5cm}

\begin{minipage}[t][1.7cm][t]{\textwidth}%
  \vspace{0pt}\raggedright\thesistitlefontnormalsize\sffamily
  %
  an der

  H\"{o}heren technischen Bundeslehr- und Versuchsanstalt Graz-Gösting
\end{minipage}

\begin{minipage}[t][2cm][t]{\textwidth}%
  \vspace{0pt}\sffamily\thesistitlefontnormalsize\raggedright
  %
  Betreuung

  \tuinfthesisbetreuung: \tuinfthesisbetreins
\end{minipage}

% we want a german date, then switch back to english
\selectlanguage{ngerman}
\begin{minipage}[t][2cm][t]{\textwidth}%
  \vspace{0pt}\sffamily\thesistitlefontnormalsize
  \begin{tabbing}%
    \hspace{44mm} \= \hspace{66mm} \= \hspace{51mm} \kill
    Graz, \tuinfthesisdate \> {\raggedright\rule{51mm}{0.5pt}} \> {\raggedright\rule{51mm}{0.5pt}} \\
    \> \begin{minipage}[t][0.5cm][t]{51mm}\centering (Unterschrift \tuinfthesisverfassung)\end{minipage}
    \> \begin{minipage}[t][0.5cm][t]{51mm}\centering (Unterschrift \tuinfthesisbetreuung)\end{minipage}
    \end{tabbing}
\end{minipage}


\selectlanguage{ngerman}
\begin{minipage}[t][2cm][t]{\textwidth}%
  \vspace{0pt}\sffamily\thesistitlefontnormalsize
  \begin{tabbing}%
    \hspace{44mm} \= \hspace{66mm} \= \hspace{51mm} \kill
    Graz, \tuinfthesisdate \> {\raggedright\rule{51mm}{0.5pt}} \> {\raggedright\rule{51mm}{0.5pt}} \\
    \> \begin{minipage}[t][0.5cm][t]{51mm}\centering (Unterschrift \tuinfthesisverfassung)\end{minipage}
    \> \begin{minipage}[t][0.5cm][t]{51mm}\centering (Unterschrift \tuinfthesisbetreuung)\end{minipage}
    \end{tabbing}
\end{minipage}


%\selectlanguage{ngerman}

\end{center}

% we want an empty page right after first titlepage
\cleardoublepage

% we're done with the titlepages, proceed with default pagenumbering
\pagenumbering{roman}

% restore baselineskip
\setlength{\baselineskip}{\tmpbaselineskip}
\setlength{\parindent}{\tmpparindent}

% back to normal geometry
\restoregeometry


%%% Local Variables:
%%% TeX-PDF-mode: t
%%% TeX-debug-bad-boxes: t
%%% TeX-parse-self: t
%%% TeX-auto-save: t
%%% reftex-plug-into-AUCTeX: t
%%% End:
 
%
% declaration
%
\textbf{
\begin{center}
Erklärung zur Verfassung der Arbeit
\end{center}
}
Hiermit erkläre ich \dots geboren am \dots in \dots, wohnhaft in \dots, dass ich diese Arbeit selbständig verfasst habe, dass ich die verwendeten Quellen und Hilfsmittel vollständig angegeben habe und dass ich die Stellen der Arbeit - einschließlich Tabellen, Karten und Abbildungen -, die anderen Werken oder dem Internet im Wortlaut oder dem Sinn nach entnommen sind, auf jeden Fall unter Angabe der Quelle als Entlehnung kenntlich gemacht habe.
\\
\\
%
%\begin{minipage}[t][2cm][t]{\textwidth}%
%  \vspace{0pt}\sffamily\thesistitlefontnormalsize
%  \begin{tabbing}%
%    \hspace{42mm} \= \hspace{66mm} \= \hspace{51mm} \kill
%    Wien, \tuinfthesisdate \> {\raggedright\rule{51mm}{0.5pt}} \> {} \\
%    \> \begin{minipage}[t][0.5cm][t]{51mm}\ (Unterschrift \tuinfthesisverfassung)\end{minipage}
%		\> \begin{minipage}[t][0.5cm][t]{51mm}\end{minipage}
%    \end{tabbing}
%\end{minipage}

\vspace{0,5 cm} 
\begin{tabular}{p{7cm}p{.5cm}l}
Graz, den {\tuinfthesisdate}
\end{tabular}% 


% Hier kommen die Unterschriten hin
\vspace{0,5 cm} 
\begin{tabular}{p{5cm}p{.5cm}l}
\dotfill \\ 
Unterschrift \tuinfthesisverfassung 
\end{tabular}% 

\vspace{4 cm}


\textbf{
\begin{center}
Erklärung zur Verfassung der Arbeit
\end{center}
}
Hiermit erkläre ich \dots geboren am \dots in \dots, wohnhaft in \dots, dass ich die verwendeten Quellen und Hilfsmittel vollständig angegeben habe und dass ich die Stellen der Arbeit - einschließlich Tabellen, Karten und Abbildungen -, die anderen Werken oder dem Internet im Wortlaut oder dem Sinn nach entnommen sind, auf jeden Fall unter Angabe der Quelle als Entlehnung kenntlich gemacht habe.
\\
\\
%
%\begin{minipage}[t][2cm][t]{\textwidth}%
%  \vspace{0pt}\sffamily\thesistitlefontnormalsize
%  \begin{tabbing}%
%    \hspace{42mm} \= \hspace{66mm} \= \hspace{51mm} \kill
%    Wien, \tuinfthesisdate \> {\raggedright\rule{51mm}{0.5pt}} \> {} \\
%    \> \begin{minipage}[t][0.5cm][t]{51mm}\ (Unterschrift \tuinfthesisverfassung)\end{minipage}
%		\> \begin{minipage}[t][0.5cm][t]{51mm}\end{minipage}
%    \end{tabbing}
%\end{minipage}

\vspace{0,5 cm} 
\begin{tabular}{p{7cm}p{.5cm}l}
Graz, den {\tuinfthesisdate}
\end{tabular}% 


% Hier kommen die Unterschriten hin
\vspace{0,5 cm} 
\begin{tabular}{p{5cm}p{.5cm}l}
\dotfill \\ 
Unterschrift \tuinfthesisverfassung 
\end{tabular}% 



%
% abstract
%
\abstractintoc
%\setlength{\parindent}{_neuer_Einzug}
%\setlength{\parskip}{12pt}

\renewcommand{\baselinestretch}{1.2}
\renewcommand{\abstractnamefont}{\normalfont\Large\bfseries}
\renewcommand{\abstractname}{Kurzfassung}
\renewcommand{\abstracttextfont}{\normalsize\nonzeroparskip}
\begin{abstract}
\lipsum[2-4]
\end{abstract}
\newpage
\renewcommand{\abstractname}{Abstract}

\begin{abstract}
\lipsum[2-4]
\end{abstract}
\newpage

\cleardoublepage

\makeatletter
\@openrighttrue % open on right pages
\makeatother


%
% table of contents, figures and tables
%
\tableofcontents
%
% for the list of XXX, we don't want pages to open right
%\makeatletter
%\@openrightfalse % open on next page
%\makeatother
%
%\listoffigures
%
%\listoftables
%
%\listofalgorithms
%


%
% start of the thesis
%
\mainmatter
\nonzeroparskip

%----------------- Kapitel Einleitung
\chapter{Einleitung}
\label{cap:einleitung}



\section{Problemstellung}

Beispiel für \cite{Akgun07}.. 

\lipsum[2-4]

\section{Ziel der Arbeit}
\lipsum[2-4]




\section{Aufbau}
Die Arbeit ist in Kapitel unterteilt, welche so weit als möglich in sich selbst abgeschlossen sind. 

\begin{description}
\item[Kapitel \ref{cap:einleitung}] gibt einen Überblick über die Problemstellung und das Ziel der Arbeit. 
\item[Kapitel \ref{cap:theorie}] führt den Leser in das Themengebiet der Softwareentwicklung und erlernen von Sprachen ein. Das Wissen über verschiedene Verfahren zum erstellen von Feedback in Sprachen soll aufbauend vermittelt werden.
\item[Kapitel \ref{cap:empirisch}] befasst sich mit der konkreten Umsetzung eines/oder mehrerer Unterrichtskonzeptes aufbauend auf den vorher elaborierten Theorieteil. 
\item[Kapitel \ref{cap:conclusion}] dient als eine Zusammenfassung der Arbeit und gibt Ausblick auf mögliche weitere Forschungsfelder.
\end{description}

%
%----------------- Kapitel Theoretischer Teil
\chapter{Theoretischer Teil}
\label{cap:theorie}

\section{Abschnitt A}
\label{cap:theorie:defA}

\section{Abschnitt B}
\label{cap:theorie:def}

\subsection{Unterabschnitt B1}
\label{cap:theorie:defC}


%----------------- Kapitel Empirischer Teil
\chapter{Praktischer Teil}
\label{cap:empirisch}

\section{Abschnitt A}
\label{cap:empirisch:defA}

\section{Abschnitt B}
\label{cap:empirisch:def}

\subsection{Unterabschnitt B1}
\label{cap:empirisch:defC}

\section{Auswertung}
\label{cap:empirisch:ausw}


%----------------- Kapitel Conclusio
\chapter{Conclusio}
\label{cap:conclusion}






%Hier beginnt der Anhang!
\appendix

%----------------- Anhang
\chapter{Beweise}


%
% bibliography
%
% AUCTeX doesn't like FILE.tex _and_ FILE.bib, so we rename FILE.bib
% to FILE-bib.bib
%
\newpage
\bibliography{literatur-bib}
\nocite{*}


%
% add a glossary
%
\newpage
\printglossary
\glsaddall

\newpage

%
% add index right after the bibliography
%
\printindex

%\newpage
%\listoftodos 

\end{document}

%%% Local Variables:
%%% TeX-PDF-mode: t
%%% TeX-debug-bad-boxes: t
%%% TeX-master: t
%%% TeX-parse-self: t
%%% TeX-auto-save: t
%%% reftex-plug-into-AUCTeX: t
%%% End: